\PassOptionsToPackage{dvipsnames}{xcolor}

\documentclass[12pt]{beamer}
\usetheme{default}
\usecolortheme{crane}

\usepackage[T1]{fontenc}
\usepackage[slovak]{babel}

\usepackage{amsmath}
\usepackage{amsmath, amssymb}
\usepackage{hyperref, url}
\usepackage{graphicx}
\usepackage{array}
\usepackage{alltt}
\usepackage{tikz}
\usetikzlibrary{matrix,calc}
\usepackage[normalem]{ulem} % \sout for struck-through rows

%\setbeamersize{text margin left=1pt,text margin right=1pt}
\setbeamertemplate{footline}[frame number]
\beamertemplatenavigationsymbolsempty

% https://www.overleaf.com/learn/latex/Using_colours_in_LaTeX
\def\blue#1{\textcolor{Cerulean}{#1}}
\def\green#1{\textcolor{LimeGreen}{#1}}

% database-related stuff
\DeclareMathOperator{\join}{\bowtie}
\DeclareMathOperator{\antijoin}{\rhd}

\DeclareMathOperator{\lubi}{lubi}
\DeclareMathOperator{\capuje}{capuje}
\DeclareMathOperator{\navstivil}{navstivil}
\DeclareMathOperator{\vypil}{vypil}
\DeclareMathOperator{\answer}{answer}

% TikZ helpers for the JOIN slides (table look-alikes of the former PNG images)
\definecolor{tblhdr}{RGB}{91,155,213}
\definecolor{tbldark}{RGB}{208,222,240}
\definecolor{tbllight}{RGB}{232,238,248}
\definecolor{tblhl}{RGB}{197,224,180}
\definecolor{vennhl}{RGB}{132,185,87}
\tikzset{
  dbt/.style={matrix of nodes, anchor=north west,
    nodes={draw=none, anchor=center, font=\fontsize{7}{8}\selectfont, minimum height=0.34cm, inner ysep=0pt, inner xsep=0.3mm},
    row 1/.style={nodes={fill=tblhdr, text=white, font=\fontsize{7}{8}\selectfont\bfseries}},
    row 2/.style={nodes={fill=tbllight}}, row 3/.style={nodes={fill=tbldark}},
    row 4/.style={nodes={fill=tbllight}}, row 5/.style={nodes={fill=tbldark}},
    row 6/.style={nodes={fill=tbllight}}, row 7/.style={nodes={fill=tbldark}},
    row 8/.style={nodes={fill=tbllight}}, row 9/.style={nodes={fill=tbldark}},
    row 10/.style={nodes={fill=tbllight}}},
  tcap/.style={anchor=south west, yshift=1pt, font=\tiny\itshape, text=gray!70!black},
}
\newcommand{\dead}[1]{\textcolor{gray}{\sout{#1}}}


\title{Relačný model a SQL}
\author{Ján Mazák}
\institute{FMFI UK Bratislava}
\date{}


\begin{document}

\frame{\titlepage}

\begin{frame}{Relačný model}
Databáza pozostáva z tabuliek (\alert{relácií}).\\[3mm]

Stĺpce --- \alert{atribúty}; každý má doménu (množinu povolených hodnôt, čiže dátový typ prípadne zúžený dodatočnými obmedzeniami). Určené pomenovaním alebo pozíciou.\\[3mm]

Riadky --- \alert{záznamy} (records, rows, tuples, n-tice).
\end{frame}

\begin{frame}[fragile]{Relačný model}
\begin{verbatim}
CREATE TABLE employees (
    id INTEGER PRIMARY KEY,
    lastName TEXT NOT NULL,
    firstName VARCHAR(255),
    age INTEGER CHECK (age >= 18)
);
\end{verbatim}
\end{frame}

\begin{frame}{Relačný model}
Reláciu možno vnímať ako predikát alebo ako (multi)množinu záznamov.\\[3mm]

V bežných DBMS sa relácia chápe ako multimnožina:
\begin{itemize}
\item na poradí riadkov nezáleží;
\item riadky v tabuľke sa môžu opakovať.
\end{itemize}

Odstránenie duplikátnych záznamov z výsledku dotazu:
\begin{alltt}
    SELECT \alert{DISTINCT} x, y FROM ...
\end{alltt}
\end{frame}

\begin{frame}[fragile]{NULL}
\begin{itemize}
\item Špeciálna hodnota \alert{NULL} zodpovedá neznámej hodnote.
\item Trojhodnotová logika, napr. \verb|NULL OR FALSE| je UNKNOWN (ani TRUE, ani FALSE).
\item Test, či je hodnota NULL: \verb|x IS NULL| / \verb|x IS NOT NULL|
\item Test, či je logický výraz UNKNOWN: \verb|(x > 0) IS UNKNOWN| (vhodné len pre debugging)
\item Pri vytváraní tabuľky možno NULL zakázať.
Inak treba starostlivo zvažovať, ako ovplyvní operátory a agregačné funkcie (napr. priemer hodnôt v stĺpci).
Nehádajte, použite dokumentáciu a testujte.
\end{itemize}
\end{frame}

\begin{frame}[fragile]{Dotazy v SQL}
Základná štruktúra:
\begin{verbatim}
SELECT attribute1 AS a1, attribute2 AS a2
FROM some_table AS t
WHERE t.attribute2 > 10
ORDER BY a1, a2
\end{verbatim}
(Všimnite si, že aliasy a1, a2 nemožno použiť vo WHERE. Kľúčové slovo AS je takmer všade nepovinné.)
\end{frame}

\begin{frame}[fragile]{Dotazy v SQL}
\begin{itemize}
\item Pri kľúčových slovách jazyka SQL sa nerozlišujú malé a veľké písmená,
ale pre dáta uložené v db áno (ak to nezmeníme napr. použitím ILIKE v podmienke za WHERE).
\item Case-sensitivity názvov tabuliek a atribútov závisí od DBMS a operačného systému.
\item Úvodzovky pre stĺpce: \verb|"atribút s medzerou v názve"|
\item Apostrofy pre konštantné reťazce: \verb|'reťazec'|
\end{itemize}

Bežné konvencie:
\begin{itemize}
\item názvy tabuliek aj atribútov lower case
\item kľúčové slová SQL upper case
\end{itemize}
\end{frame}

\begin{frame}[fragile]{Dotazy v SQL}
Vo výsledku nemusia byť len pôvodné hodnoty atribútov, ale aj čosi z nich vyrátané (napr. aritmetické výrazy zložené z konštánt, funkcií implementovaných v db a hodnôt atribútov daného riadka).\\[3mm]

\begin{verbatim}
SELECT
    concat(e.firstname,' ',e.lastname) AS ename,
    0.8 * e.salary AS salaryAfterTax,
    (CASE
        WHEN e.bonus IS NULL THEN e.salary
        ELSE e.bonus + e.salary
    END) AS total_salary
FROM employee AS e
WHERE dept_id >= 20 AND lower(e.firstname) = 'john'
ORDER BY 0.8 * e.salary
\end{verbatim}
\end{frame}

\begin{frame}[fragile]{Zoradenie záznamov}
Zoznam zamestnancov usporiadaný od najvyššieho platu po najnižší, pri rovnakom plate abecedne:
\begin{alltt}
SELECT name, salary, department
FROM employee
\alert{ORDER BY} salary DESC, name ASC
\end{alltt}
\end{frame}

\begin{frame}{Join}
\alert{Join} --- spojenie záznamov z dvoch tabuliek.\\[3mm]

(INNER) JOIN je podmnožina karteziánskeho súčinu tabuliek (každý riadok s každým, dvojicu záznamov spojíme do jedného dlhšieho)
špecifikovaná dodatočnými podmienkami na prepájanie.
\end{frame}

\begin{frame}[fragile]{Join --- karteziánsky súčin}
\hspace*{-0.3cm}\begin{tikzpicture}[font=\fontsize{7}{8}\selectfont]
  \matrix (emp) [dbt, anchor=north west, column 1/.style={nodes={minimum width=0.95cm}}, column 2/.style={nodes={minimum width=1.15cm}}] at (0,0) {
    name & dept\_id \\
    John & 10 \\
    Thomas & 20 \\
    Joe & 40 \\
  };
  \node[tcap] at (emp.north west) {employee};
  \matrix (dep) [dbt, anchor=north west, column 1/.style={nodes={minimum width=1.15cm}}, column 2/.style={nodes={minimum width=1.50cm}}] at ($(emp.north east)+(0.3,0)$) {
    dept\_id & name \\
    10 & Accounting \\
    20 & PR \\
    30 & Development \\
  };
  \node[tcap] at (dep.north west) {department};
  \node (conn) at ($(emp.east)!0.5!(dep.west)$) {$\times$};
  \matrix (res) [dbt, anchor=north west, column 1/.style={nodes={minimum width=0.95cm}}, column 2/.style={nodes={minimum width=1.20cm}}, column 3/.style={nodes={minimum width=1.20cm}}, column 4/.style={nodes={minimum width=1.50cm}}] at ($(dep.north east)+(0.5,0)$) {
    e.name & e.dept\_id & d.dept\_id & d.name \\
    John & 10 & 10 & Accounting \\
    John & 10 & 20 & PR \\
    John & 10 & 30 & Development \\
    Thomas & 20 & 10 & Accounting \\
    Thomas & 20 & 20 & PR \\
    Thomas & 20 & 30 & Development \\
    Joe & 40 & 10 & Accounting \\
    Joe & 40 & 20 & PR \\
    Joe & 40 & 30 & Development \\
  };
  \node at ($(res.west |- conn)+(-0.25,0)$) {$=$};
\end{tikzpicture}

\begin{alltt}
SELECT * FROM employee, department
\end{alltt}
\end{frame}

\begin{frame}[fragile]{Join --- INNER JOIN (= JOIN)}
\hspace*{-0.3cm}\begin{tikzpicture}[font=\fontsize{7}{8}\selectfont]
  \matrix (emp) [dbt, anchor=north west, column 1/.style={nodes={minimum width=0.95cm}}, column 2/.style={nodes={minimum width=1.15cm}}] at (0,0) {
    name & dept\_id \\
    John & 10 \\
    Thomas & 20 \\
    Joe & 40 \\
  };
  \node[tcap] at (emp.north west) {employee};
  \matrix (dep) [dbt, anchor=north west, column 1/.style={nodes={minimum width=1.15cm}}, column 2/.style={nodes={minimum width=1.50cm}}] at ($(emp.north east)+(0.5,0)$) {
    dept\_id & name \\
    10 & Accounting \\
    20 & PR \\
    30 & Development \\
  };
  \node[tcap] at (dep.north west) {department};
  \node (conn) [font=\scriptsize] at ($(emp.east)!0.5!(dep.west)$) {JOIN};
  \matrix (res) [dbt, anchor=north west, column 1/.style={nodes={minimum width=0.95cm}}, column 2/.style={nodes={minimum width=1.20cm}}, column 3/.style={nodes={minimum width=1.20cm}}, column 4/.style={nodes={minimum width=1.50cm}}] at ($(dep.north east)+(0.5,0)$) {
    e.name & e.dept\_id & d.dept\_id & d.name \\
    John & 10 & 10 & Accounting \\
    \dead{John} & \dead{10} & \dead{20} & \dead{PR} \\
    \dead{John} & \dead{10} & \dead{30} & \dead{Development} \\
    \dead{Thomas} & \dead{20} & \dead{10} & \dead{Accounting} \\
    Thomas & 20 & 20 & PR \\
    \dead{Thomas} & \dead{20} & \dead{30} & \dead{Development} \\
    \dead{Joe} & \dead{40} & \dead{10} & \dead{Accounting} \\
    \dead{Joe} & \dead{40} & \dead{20} & \dead{PR} \\
    \dead{Joe} & \dead{40} & \dead{30} & \dead{Development} \\
  };
  \node at ($(res.west |- conn)+(-0.25,0)$) {$=$};
  \begin{scope}[shift={(1.2,-3.2)}, scale=0.75]
    \begin{scope}
      \clip (-0.6,0) ellipse (1.05cm and 0.5cm);
      \fill[vennhl] (0.6,0) ellipse (1.05cm and 0.5cm);
    \end{scope}
    \draw (-0.6,0) ellipse (1.05cm and 0.5cm);
    \draw (0.6,0) ellipse (1.05cm and 0.5cm);
    \node[font=\tiny] at (-0.9,0) {table1};
    \node[font=\tiny] at (0.9,0) {table2};
    \node[font=\scriptsize] at (0,0.8) {INNER JOIN};
  \end{scope}
\end{tikzpicture}
\vspace{-0.2em}\rule{\linewidth}{0.3pt}\vspace{-0.2em}

\footnotesize
\begin{alltt}
SELECT * FROM employee e, department d
WHERE e.dept\_id = d.dept\_id
\end{alltt}
\vspace{-0.4em}\rule{\linewidth}{0.3pt}\vspace{-0.2em}
\begin{alltt}
SELECT * FROM employee e
    JOIN department d ON e.dept\_id = d.dept\_id
\end{alltt}
\vspace{-0.4em}\rule{\linewidth}{0.3pt}\vspace{-0.2em}
\textbf{Pozor:} \texttt{NATURAL JOIN} spája na všetky stĺpce s rovnakým názvom, v tomto príklade teda na \texttt{name} aj \texttt{dept\_id}. Výsledok je prázdny (0 riadkov), kým \texttt{JOIN ... ON} vráti 3 riadky.
\end{frame}

\begin{frame}[fragile]{Join --- LEFT [OUTER] JOIN}
\hspace*{-0.3cm}\begin{tikzpicture}[font=\fontsize{7}{8}\selectfont]
  \matrix (emp) [dbt, anchor=north west, column 1/.style={nodes={minimum width=0.95cm}}, column 2/.style={nodes={minimum width=1.15cm}}] at (0,0) {
    name & dept\_id \\
    John & 10 \\
    Thomas & 20 \\
    Joe & 40 \\
  };
  \node[tcap] at (emp.north west) {employee};
  \matrix (dep) [dbt, anchor=north west, column 1/.style={nodes={minimum width=1.15cm}}, column 2/.style={nodes={minimum width=1.50cm}}, row 5/.style={nodes={fill=tblhl}}] at ($(emp.north east)+(0.8,0)$) {
    dept\_id & name \\
    10 & Accounting \\
    20 & PR \\
    30 & Development \\
    10 & Human res. \\
  };
  \node[tcap] at (dep.north west) {department};
  \node (conn) [align=center, font=\scriptsize] at ($(emp.east)!0.5!(dep.west)$) {LEFT\\JOIN};
  \matrix (res) [dbt, anchor=north west, column 1/.style={nodes={minimum width=0.95cm}}, column 2/.style={nodes={minimum width=1.20cm}}, column 3/.style={nodes={minimum width=1.20cm}}, column 4/.style={nodes={minimum width=1.50cm}}, row 3/.style={nodes={fill=tblhl}}] at ($(dep.north east)+(0.5,0)$) {
    e.name & e.dept\_id & d.dept\_id & d.name \\
    John & 10 & 10 & Accounting \\
    John & 10 & 10 & Human res. \\
    Thomas & 20 & 20 & PR \\
    Joe & 40 & null & null \\
  };
  \node at ($(res.west |- conn)+(-0.25,0)$) {$=$};
\end{tikzpicture}
\par\vspace{1.0cm}\noindent
\begin{minipage}[c]{0.56\linewidth}

\begin{alltt}
SELECT *
FROM employee as e
    LEFT JOIN department as d
    ON e.dept\_id = d.dept\_id
\end{alltt}
\end{minipage}%
\begin{minipage}[c]{0.40\linewidth}\centering
\begin{tikzpicture}[font=\fontsize{7}{8}\selectfont]
  \begin{scope}
    \fill[vennhl] (-0.6,0) ellipse (1.05cm and 0.5cm);
    \draw (-0.6,0) ellipse (1.05cm and 0.5cm);
    \draw (0.6,0) ellipse (1.05cm and 0.5cm);
    \node[font=\tiny] at (-0.9,0) {table1};
    \node[font=\tiny] at (0.9,0) {table2};
    \node[font=\scriptsize] at (0,0.8) {LEFT JOIN};
  \end{scope}
\end{tikzpicture}
\end{minipage}
\end{frame}

\begin{frame}[fragile]{Join --- RIGHT [OUTER] JOIN}
\hspace*{-0.5cm}\begin{tikzpicture}[font=\fontsize{7}{8}\selectfont]
  \matrix (dep) [dbt, anchor=north west, column 1/.style={nodes={minimum width=1.15cm}}, column 2/.style={nodes={minimum width=1.50cm}}, row 5/.style={nodes={fill=tblhl}}] at (0,0) {
    dept\_id & name \\
    10 & Accounting \\
    20 & PR \\
    30 & Development \\
    10 & Human res. \\
  };
  \node[tcap] at (dep.north west) {department};
  \matrix (emp) [dbt, anchor=north west, column 1/.style={nodes={minimum width=0.95cm}}, column 2/.style={nodes={minimum width=1.15cm}}] at ($(dep.north east)+(0.8,0)$) {
    name & dept\_id \\
    John & 10 \\
    Thomas & 20 \\
    Joe & 40 \\
  };
  \node[tcap] at (emp.north west) {employee};
  \node (conn) [align=center, font=\scriptsize] at ($(dep.east)!0.5!(emp.west)$) {RIGHT\\JOIN};
  \matrix (res) [dbt, anchor=north west, column 1/.style={nodes={minimum width=1.20cm}}, column 2/.style={nodes={minimum width=1.50cm}}, column 3/.style={nodes={minimum width=0.95cm}}, column 4/.style={nodes={minimum width=1.20cm}}, row 3/.style={nodes={fill=tblhl}}] at ($(emp.north east)+(0.3,0)$) {
    d.dept\_id & d.name & e.name & e.dept\_id \\
    10 & Accounting & John & 10 \\
    10 & Human res. & John & 10 \\
    20 & PR & Thomas & 20 \\
    null & null & Joe & 40 \\
  };
  \node at ($(res.west |- conn)+(-0.25,0)$) {$=$};
\end{tikzpicture}
\par\vspace{1.0cm}\noindent
\begin{minipage}[c]{0.56\linewidth}

\begin{alltt}
SELECT *
FROM department as d
    RIGHT JOIN employee as e
    ON e.dept\_id = d.dept\_id
\end{alltt}
\end{minipage}%
\begin{minipage}[c]{0.40\linewidth}\centering
\begin{tikzpicture}[font=\fontsize{7}{8}\selectfont]
  \begin{scope}
    \fill[vennhl] (0.6,0) ellipse (1.05cm and 0.5cm);
    \draw (-0.6,0) ellipse (1.05cm and 0.5cm);
    \draw (0.6,0) ellipse (1.05cm and 0.5cm);
    \node[font=\tiny] at (-0.9,0) {table1};
    \node[font=\tiny] at (0.9,0) {table2};
    \node[font=\scriptsize] at (0,0.8) {RIGHT JOIN};
  \end{scope}
\end{tikzpicture}
\end{minipage}
\end{frame}

\begin{frame}[fragile]{Join --- FULL [OUTER] JOIN}
\hspace*{-0.3cm}\begin{tikzpicture}[font=\fontsize{7}{8}\selectfont]
  \matrix (emp) [dbt, anchor=north west, column 1/.style={nodes={minimum width=0.95cm}}, column 2/.style={nodes={minimum width=1.15cm}}] at (0,0) {
    name & dept\_id \\
    John & 10 \\
    Joe & 40 \\
  };
  \node[tcap] at (emp.north west) {employee};
  \matrix (dep) [dbt, anchor=north west, column 1/.style={nodes={minimum width=1.15cm}}, column 2/.style={nodes={minimum width=1.50cm}}] at ($(emp.north east)+(0.8,0)$) {
    dept\_id & name \\
    10 & Accounting \\
    30 & Development \\
  };
  \node[tcap] at (dep.north west) {department};
  \node (conn) [align=center, font=\scriptsize] at ($(emp.east)!0.5!(dep.west)$) {FULL\\JOIN};
  \matrix (res) [dbt, anchor=north west, column 1/.style={nodes={minimum width=0.95cm}}, column 2/.style={nodes={minimum width=1.20cm}}, column 3/.style={nodes={minimum width=1.20cm}}, column 4/.style={nodes={minimum width=1.50cm}}] at ($(dep.north east)+(0.5,0)$) {
    e.name & e.dept\_id & d.dept\_id & d.name \\
    John & 10 & 10 & Accounting \\
    Joe & 40 & null & null \\
    null & null & 30 & Development \\
  };
  \node at ($(res.west |- conn)+(-0.25,0)$) {$=$};
\end{tikzpicture}
\par\vspace{1.0cm}\noindent
\begin{minipage}[c]{0.56\linewidth}

\begin{alltt}
SELECT *
FROM employee e
    FULL JOIN department d
    ON e.dept\_id = d.dept\_id
\end{alltt}
\end{minipage}%
\begin{minipage}[c]{0.40\linewidth}\centering
\begin{tikzpicture}[font=\fontsize{7}{8}\selectfont]
  \begin{scope}
    \fill[vennhl] (-0.6,0) ellipse (1.05cm and 0.5cm);
    \fill[vennhl] (0.6,0) ellipse (1.05cm and 0.5cm);
    \draw (-0.6,0) ellipse (1.05cm and 0.5cm);
    \draw (0.6,0) ellipse (1.05cm and 0.5cm);
    \node[font=\tiny] at (-0.9,0) {table1};
    \node[font=\tiny] at (0.9,0) {table2};
    \node[font=\scriptsize] at (0,0.8) {FULL JOIN};
  \end{scope}
\end{tikzpicture}
\end{minipage}
\end{frame}

\begin{frame}[fragile]{Negácia}
Negáciu možno vyjadriť pomocou \alert{NOT EXISTS}:
\begin{alltt}
/* zamestnanci, ktorí nemajú podriadených */
SELECT emp.name FROM employee \blue{emp}
WHERE \alert{NOT EXISTS} (
    SELECT 1 FROM employee \green{emp2}
    WHERE \green{emp2}.superior = \blue{emp}.emp_id
)
\end{alltt}
(Nezáleží na tom, ktoré stĺpce sú vymenované za SELECT vo vnútornom dotaze, pretože EXISTS len testuje, či sa tam nachádza aspoň jeden riadok.)
\end{frame}

\begin{frame}[fragile]{Negácia}
Samotný operátor EXISTS bez NOT je v podstate JOIN.
\begin{alltt}
/* zamestnanci, ktorí majú podriadených */
SELECT DISTINCT emp.name
FROM employee \blue{emp}
WHERE \alert{EXISTS} (
    SELECT 1 FROM employee \green{emp2}
    WHERE \green{emp2}.superior = \blue{emp}.emp_id
)

SELECT DISTINCT emp.name
FROM employee \blue{emp}, employee \green{emp2}
WHERE \green{emp2}.superior = \blue{emp}.emp_id
\end{alltt}
\end{frame}

\begin{frame}[fragile]{Všeobecný kvantifikátor}
Jazyk SQL nemá prostriedky na priame vyjadrenie všeobecného kvantifikátora.
Postupovať možno v duchu
$$
    \forall x\ P(x)\qquad \Leftrightarrow\qquad \lnot\ \exists x\ \lnot\, P(x).
$$
Zamestnanci, ktorí majú najvyšší plat spomedzi \alert{všetkých} $\Leftrightarrow$ zamestnanci, ku ktorým \alert{neexistuje} zamestnanec s vyšším platom.
\begin{alltt}
SELECT e.name
FROM employee e
WHERE \blue{NOT EXISTS} (SELECT 1
                  FROM employee e2
                  WHERE e2.salary > e.salary)
\end{alltt}
\end{frame}

\begin{frame}[fragile]{Vnorené dotazy (subqueries)}
Okrem operátora EXISTS sa možno pomocou operátora IN (resp. NOT IN) pýtať na prítomnosť v množine.
\begin{alltt}
SELECT name
FROM employee e
WHERE e.dept_id \alert{IN} (SELECT dept_id
                   FROM department d
                   WHERE d.location = 'New York')
\end{alltt}
NOT IN sa optimalizuje horšie ako antijoin (NOT EXISTS), lebo kvôli NULL ho nemožno jednoducho prepísať na antijoin.
NOT IN má tiež nevýhodu, ak sa môže vyskytnúť NULL.
\end{frame}

\begin{frame}[fragile]{Vnorené dotazy (subqueries)}
Ak je výsledkom vnoreného dotazu tabuľka 1x1, možno ju použiť ako skalár.
\begin{alltt}
SELECT name
FROM employee e
WHERE \alert{e.dept_id =} \blue{(SELECT dept_id
                  FROM department d
                  WHERE d.location = 'New York')}
\end{alltt}
Ak by vo vnorenom dotaze vyšlo viac riadkov, dôjde k chybe; ak žiaden, výsledok je NULL.\\
(Takýmto chybám treba čo najviac predchádzať, napr. použitím UNIQUE už pri vytváraní tabuľky.)
\end{frame}

\begin{frame}{Množinové operácie}
Bez zachovania násobnosti riadkov (matematické množiny):
\begin{itemize}
\item UNION --- zjednotenie
\item INTERSECT --- prienik
\item EXCEPT --- rozdiel množín
\end{itemize}
So zachovaním násobnosti riadkov (multimnožiny):
\begin{itemize}
\item UNION ALL --- zjednotenie
\item INTERSECT ALL --- prienik
\item EXCEPT ALL --- rozdiel množín
\end{itemize}
\end{frame}

\begin{frame}[fragile]{Množinové operácie}
\begin{minipage}{.4\pdfpagewidth}
\begin{alltt}
SELECT nazov
FROM kry

\alert{UNION}

SELECT nazov
FROM stromy
\end{alltt}
\end{minipage}
\begin{minipage}{.4\pdfpagewidth}
\begin{alltt}
SELECT nazov
FROM dreviny

\alert{EXCEPT}

SELECT nazov
FROM stromy
\end{alltt}
\end{minipage}
\end{frame}

\begin{frame}[fragile]{Mačky nežerúce myši --- bežná chyba}
Mačky možno rozdeliť do 4 disjunktných skupín:
\begin{itemize}
\item[A.] nežerie nič
\item[B.] žerie niečo, ale nie myši
\item[C.] žerie myši, ale nič iné
\item[D.] žerie myši aj niečo iné
\end{itemize}
$\hbox{mačky nežerúce myši} = A\cup B$
$\hbox{mačky žerúce iné ako myši} = B\cup D$
\\[3mm]

Ak je dotaz o \uv{mačkách nežerúcich myši} a máme za WHERE \uv{$\hbox{potrava} <> \hbox{myši}$}, asi je to nesprávne.
Pre určenie toho, či mačka žerie myši, je úplne irelevantná existencia záznamov o tom, že žerie niečo iné ako myši.

(Dve množiny popísané vyššie však budú totožné, ak každá mačka žerie práve jednu potravu.)
\end{frame}

\begin{frame}[fragile]{CREATE TEMPORARY TABLE}
Používateľ si vie sám vytvoriť dočasnú tabuľku, ktorá zanikne pri ukončení \emph{session} (odpojení od db). Voľby ON COMMIT DROP alebo ON COMMIT DELETE ROWS ho naopak viažu na koniec transakcie.
(Výsledok dotazu sa vypočíta a uloží pri vytvorení tabuľky; pri ďalšom použití ho nepočítame nanovo.)
\begin{alltt}
\alert{CREATE TEMPORARY TABLE} \blue{employee_details} AS (
    SELECT e.emp_id, e.name, e.dept_id,
           d.name AS dept_name, d.location
    FROM employee e, department d
    WHERE e.dept_id = d.dept_id
);
SELECT * FROM \blue{employee_details} WHERE ...;
\end{alltt}
Hodí sa to napr. ak chceme opakovane využívať ten istý komplikovaný join.
\end{frame}


\begin{frame}[fragile]{VIEW}
Prístup k dátam možno uľahčiť vytvorením pohľadu (VIEW).
\begin{alltt}
\alert{CREATE VIEW} \blue{employee_details} AS (
    SELECT e.emp_id, e.name, e.dept_id,
           d.name AS dept_name, d.location
    FROM employee e, department d
    WHERE e.dept_id = d.dept_id
);

SELECT * FROM \blue{employee_details} WHERE ...;
\end{alltt}
VIEW je permanentný objekt v db. Jeho riadky však nemusia byť materializované.
Cez VIEW možno meniť dáta len za prísnych podmienok (podľa štandardu SQL VIEW vyberá z jednej tabuľky, bez JOIN-u, DISTINCT, GROUP BY, agregácií a UNION). VIEW s JOIN-om preto nemožno priamo meniť.
Vkladať dáta do VIEW nikdy nie je dobrý nápad.
\end{frame}


\begin{frame}[fragile]{Funkcie pre jednotlivé dátové typy}
Pri praktickej práci v SQL využívame rôzne funkcie pre špecifické dátové typy.
\begin{itemize}
\item dátum a čas
\item reťazce
\item numerické výpočty
\item pretypovanie
\item \dots
\item často špecifické pre konkrétny DBMS
\item nebudeme s nimi strácať čas, ľahko sa rieši za pochodu (dokumentácia / LLM)
\end{itemize}
\end{frame}


\begin{frame}[fragile]{How to write SQL}
Zapisujte dotazy čitateľne a zrozumiteľne.\\
Vyhýbajte sa zápisom vedúcim k neefektívnemu výpočtu.\\[3mm]
\begin{itemize}
\item {\scriptsize\url{https://github.com/jarulraj/sqlcheck/blob/master/README.md#query-anti-patterns}}
\item {\scriptsize\url{https://www.stratascratch.com/blog/best-practices-to-write-sql-queries-how-to-structure-your-code/}}
\item {\scriptsize\url{https://www.sqlstyle.guide/}}
\end{itemize}
\vspace*{3mm}
Prečítajte si tieto konvencie \emph{teraz} a nezaťažujte budúcich kolegov nevhodnými zápismi.
(Berte ich však ako odporúčania; rôzne zdroje v niektorých veciach nesúhlasia. Zvyčajne sa pri existujúcom kóde riadime miestnymi konvenciami.)
\end{frame}


\begin{frame}{Literatúra}
\begin{itemize}
\item {\scriptsize\url{https://cs186berkeley.net/notes/note1/}}
\item {\scriptsize\url{https://cs186berkeley.net/notes/note2/}}
\item {\scriptsize\url{https://www.postgresqltutorial.com/} (Sections 1, 2, 3)}
\item {\scriptsize\url{https://www.w3schools.com/sql/sql_intro.asp}}
\item {\scriptsize\url{https://sqlbolt.com/lesson/select_queries_introduction}}
\item {\scriptsize\url{https://www.learnsqlonline.org/}}
\item {\scriptsize\url{https://www.postgresqltutorial.com/postgresql-views/managing-postgresql-views/}}
\item {\scriptsize\url{https://drive.google.com/file/d/1HCq2KMZ05UvtXGe1nTqNmkwhc3X-NLI3/view}}
\end{itemize}
\end{frame}


\begin{frame}{Na zamyslenie}
\begin{itemize}
\item Akú časovú zložitosť má výpočet DISTINCT (v závislosti od počtu riadkov za predpokladu, že sa všetky riadky zmestia do RAM)?
\item Join tabuliek sa dá počítať len z dvoch tabuliek naraz. Ak ich máme viac, join rátame postupne
    (predstavte si operátorový strom pre binárnu operáciu, v listoch sú pôvodné tabuľky a v koreni výsledok joinu).
    Koľko je možných postupov na výpočet joinu $4$ tabuliek? Odhadnite počet postupov pre $n$ tabuliek.
\item Vymyslite postup výpočtu INTERSECT. Akú má časovú zložitosť v závislosti od počtu riadkov relácií na vstupe? (Chcete lepšiu ako kvadratickú zložitosť.)
\end{itemize}
\end{frame}


\begin{frame}{Úlohy: SQL}
Databáza: \emph{osoba}(meno), \emph{pozna}(kto, koho)
\begin{itemize}
    \item osoby, ktoré poznajú sysľa
    \item osoby, ktoré poznajú aspoň dve entity\\ (nemusia to byť osoby)
    \item osoby, ktoré nepoznajú nič a nikoho
    \item osoby, ktoré nepoznajú žiadne iné osoby
    \item osoby, ktoré poznajú iba Jožka
    \item osoby, ktoré pozná presne jedna osoba
    \item osoby, ktoré poznajú všetkých známych svojich známych
%    \item osoby, ktoré majú všetky vzťahy symetrické
\end{itemize}
\end{frame}


\end{document}


